% SINT Protocol — SPAI 2026 Extended Abstract
% 1st IJCAI Workshop on Safe Physical AI (SPAI 2026)
% Venue: IJCAI/ECAI 2026, Bremen, Germany, August 15-17, 2026
% Submit at: https://openreview.net/group?id=ijcai.org/IJCAI-ECAI/2026/Workshop/SPAI

\documentclass{article}

\pdfpagewidth=8.5in
\pdfpageheight=11in

% --- IJCAI 2026 style ---
% Download ijcai26.sty from: https://www.ijcai.org/authors_kit
% Then uncomment the next line and remove the \usepackage{geometry} line below.
% \usepackage{ijcai26}

% Fallback geometry for local preview / arXiv preprint (remove when using ijcai26.sty):
\usepackage[margin=1in,top=0.9in,bottom=0.9in]{geometry}

\pdfinfo{
/TemplateVersion (IJCAI.2026.0)
/Title (SINT Protocol: A Runtime Authorization Framework for Physical AI Agents)
/Subject (Safe Physical AI, SPAI 2026, IJCAI/ECAI 2026)
/Keywords (physical AI, runtime authorization, capability tokens, robot safety, OWASP ASI)
}

\usepackage{times}
\usepackage{url}
\usepackage[hidelinks]{hyperref}
\usepackage[utf8]{inputenc}
\usepackage[T1]{fontenc}
\usepackage[small]{caption}
\usepackage{amsmath}
\usepackage{booktabs}
\usepackage{microtype}
\usepackage{enumitem}
\usepackage{xcolor}

\setlength{\parskip}{3pt}
\setlist[itemize]{noitemsep, topsep=2pt, parsep=0pt, partopsep=0pt}

% -----------------------------------------------------------------------

\title{\textbf{SINT Protocol: A Runtime Authorization Framework\\
for Physical AI Agents}\\[4pt]
{\normalsize \textit{Extended Abstract — 1st IJCAI Workshop on Safe Physical AI (SPAI 2026)}}}

% For blind review, use: \author{}
% For camera-ready, use the IJCAI \affiliations + \emails macros:
%
% \author{
%   Author Name$^{1}$\and
%   Author Name$^{2}$\\
%   \affiliations
%   $^{1}$Institution One, City, Country\\
%   $^{2}$Institution Two, City, Country\\
%   \emails
%   author1@example.com, author2@example.com
% }
\author{}  % blind review

\date{}

\begin{document}

\maketitle
\thispagestyle{empty}

% -----------------------------------------------------------------------

\begin{abstract}
Physical AI agents — LLM-driven systems controlling robots, drones, and actuators — introduce a
safety failure mode that neither classical robot safety (IEC~62443, ISO~10218) nor LLM alignment
research fully addresses: \textit{an authenticated, aligned agent issuing a physically irreversible
command without appropriate human oversight, at the wrong moment, in a degraded environment.}

We present \textbf{SINT Protocol}, an open-source runtime authorization framework that enforces
six safety invariants at every agent action boundary.
\end{abstract}

% -----------------------------------------------------------------------

\section*{Safety Invariants}

\begin{itemize}

\item \textbf{Capability confinement.}
Every action is authorized by a signed, expiring Ed25519 capability token scoped to a specific
resource, action set, and physical constraint envelope (max velocity, max force, geofence polygon).
Tokens cannot be self-issued or silently expanded; delegation is \emph{attenuation-only} with
cryptographic enforcement.

\item \textbf{Tier-based human oversight.}
Actions are classified into four approval tiers: \texttt{T0\_observe} (auto-allow),
\texttt{T1\_prepare} (audit), \texttt{T2\_act} (human review), \texttt{T3\_commit} (explicit
sign-off with M-of-N quorum). Classification is driven by 50+ pattern-matching rules covering
ROS~2 topics, MAVLink commands, and MCP tool calls. Shell and code-execution tool names
(\texttt{bash}, \texttt{exec}, \texttt{eval}) are explicitly classified at \texttt{T3\_commit}
to address ASI05~\cite{owasp-asi}.

\item \textbf{Behavioral drift detection (CSML).}
The Composite Safety-Model Latency score tracks an agent's out-of-policy action rate, timing
anomalies, and safety event frequency across a sliding evidence window. When CSML exceeds a
threshold~$\theta$, the agent's tier is automatically escalated — making behavioral drift a
\emph{live enforcement trigger} rather than a post-incident finding. Frontier LLMs exhibit up to
4.8$\times$ differences in out-of-policy action rates on identical tasks~\cite{rosclaw}; SINT
makes this variance enforcement-visible.

\item \textbf{Environment-adaptive constraint tightening.}
Token constraints are necessarily static at issuance. The \texttt{DynamicEnvelopePlugin} maps
real-time sensor state to tighter effective limits:
$v_{\text{eff}} = \min(\texttt{token.maxVelocityMps},\; d_{\text{obs}} \cdot r_f)$.
This closes the gap identified in ROSClaw~\cite{rosclaw} between "configurable safety envelope"
and "environment-adaptive safety envelope."

\item \textbf{Active safety plugins.}
The \texttt{GoalHijackPlugin}~(ASI01) provides 5-layer heuristic detection: prompt injection,
role override, semantic escalation, exfiltration probes, and cross-agent injection.
The \texttt{MemoryIntegrityChecker}~(ASI06) detects replay attacks, privilege claims, history
overflow, \emph{cross-session continuity injection} (hallucinated prior-auth claims),
\emph{credential read funnels} ($\geq$3 secret reads followed by a write), and
\emph{action velocity loops} (last 5 actions on an identical resource — hallucination loop
indicator). The \texttt{DefaultSupplyChainVerifier}~(ASI04) validates model fingerprint hashes
at runtime to detect tampered tools.

\item \textbf{Stop button (EU AI Act Art.~14(4)(e)).}
The \texttt{CircuitBreakerPlugin} allows operators to manually \texttt{trip()} the circuit,
instantly blocking all actions from a rogue agent. $N$ consecutive denials auto-open the circuit;
HALF\_OPEN probe recovery ensures self-healing after incident resolution.

\end{itemize}

\section*{Swarm Safety}

For multi-agent deployments, \texttt{SwarmCoordinator} enforces collective constraints before
individual gateway calls: maximum concurrent actors in T2+ tier, total kinetic energy ceiling
($\sum \tfrac{1}{2}mv^2$), minimum inter-agent distance, and maximum escalated fraction. These
collective constraints cannot be captured by per-agent token scoping alone.

\section*{Operator Interface Layer}

SINT's governance-first operator layer comprises: (i)~\texttt{@sint/memory} — a ledger-backed
operator memory system where every recall emits a typed \texttt{SintLedgerEvent};
(ii)~a voice-first SINT Command HUD routing every voice command through
\texttt{PolicyGateway.intercept()} before execution; (iii)~a \texttt{ProactiveEscalationEngine}
that monitors CSML drift across multi-turn windows and emits operator notifications at $\theta{=}0.3$;
(iv)~hierarchical token delegation (\texttt{sint\_\_delegate\_to\_agent}, max depth~3) with
atomic cascade revocation via DFS subtree traversal; and (v)~a REST API surface
(\texttt{/v1/memory}, \texttt{/v1/delegations}, \texttt{/v1/csml}).

\section*{Implementation \& Compliance}

SINT is implemented in TypeScript across \textbf{41~packages}, with \textbf{1,772~tests} covering
core, gateway, bridges, and the avatar layer. Bridge adapters exist for ROS~2, MAVLink~v2,
MCP tool calls, Google A2A, MQTT/CoAP, OPC-UA, Open-RMF, and Sparkplug~B. The Evidence Ledger
implements a SHA-256 hash chain with optional TEE attestation (Intel SGX, ARM TrustZone, AMD SEV).

SINT covers all 10 OWASP Agentic Top 10 categories (ASI01–ASI10)~\cite{owasp-asi}. Coverage is
formally declared in a machine-readable conformance mapping and regression-tested across 29
attack/safe-case fixture pairs that validate each ASI category against the gateway choke point.

\section*{Position}

Safe physical AI requires that authorization be treated as a first-class \emph{runtime} concern —
not a deployment configuration or a prompt engineering problem. The safety properties that
matter (who approved this action, was the envelope respected, was the agent's drift within bounds)
must be cryptographically auditable and enforced \emph{before} hardware execution, not after.
SINT Protocol is a concrete, open-source instantiation of this position.

{\small \textbf{Open source:} \url{https://github.com/sint-ai/sint-protocol}}

% -----------------------------------------------------------------------

% IJCAI uses named.bst (author-year). With ijcai26.sty active, use:
% \bibliographystyle{named}
% \bibliography{references}
% For self-contained fallback we use the inline thebibliography below.
\bibliographystyle{named}
\begin{thebibliography}{9}

\bibitem{rosclaw}
Cardenas, I.S., Arnett, M.A., Yeo, N.C., Sah, L., Kim, J.-H. (2026).
ROSClaw: An OpenClaw ROS~2 Framework for Agentic Robot Control and Interaction.
\textit{arXiv:2603.26997}.

\bibitem{owasp-asi}
OWASP Foundation (2025).
OWASP Agentic AI Security Top 10 — Agentic Security Initiative.
\url{https://owasp.org/www-project-top-10-for-large-language-model-applications/}.

\bibitem{iec62443}
IEC 62443-3-3:2013.
Industrial communication networks — IT security for networks and systems.

\bibitem{iso10218}
ISO 10218-1:2011.
Robots and robotic devices — Safety requirements for industrial robots.

\bibitem{nist-ai-rmf}
NIST (2023).
Artificial Intelligence Risk Management Framework (AI RMF 1.0).

\bibitem{nist-sp800-82}
NIST SP 800-82 Rev.~3 (2023).
Guide to Operational Technology (OT) Security.

\end{thebibliography}

\end{document}
